\section{Residual consequences in minimal counterexamples}\label{sec:paper-residuals}
Throughout this section, unless stated otherwise, $D$ is in the MIN setting
of Section~\ref{sec:paper-notation}. Put
\[
 \eta(v)=\mu(v)-1+d^{--}(v)-d^{++}(v),\qquad
 R=\sum_v\eta(v)=2M-n.
\]
The last identity counts exact directed distance-two pairs by their two
endpoints. Unlike the core certificate, the inequalities below use original
arc-set minimality.

\subsection{The original boundary and partner estimates}
We collect the local estimates used here. Full proofs are given in the
supplement so that the external-boundary convention is explicit.

\begin{lemma}[Original deletion estimates]\label{lem:paper-packet}
Every original gap belongs to $\{1,2\}$. If $\partial S\ne\varnothing$,
then $|\partial^2S|<|\partial S|$. Moreover $\eta(v)\ge0$ for every $v$,
and
\[
 \eta(v)=2\mu(v)+g(v)-1-|P(v)|.
\]
Fix $v$ with $P(v)\ne\varnothing$, and write
\[
 \begin{gathered}
 C=\partial P(v),\quad X=\partial^2P(v),\quad J=N^{--}(v)\setminus C,\\
 B=J\cap\mathcal M(v),\quad k=|B|,\quad s=|J\cap N^+(v)|.
 \end{gathered}
\]
Delete the $s$ arcs from $v$ into $J$, obtaining $\widehat D$. Then
\[
 N_{\widehat D}^{++}(v)\subseteq
 E=(\mathcal M(v)\setminus B)\,\dot\cup\,X,
 \qquad L=E\setminus N_{\widehat D}^{++}(v)
\]
satisfies
\begin{equation}\label{eq:paper-packet-bound}
 |L|\le\eta(v)-2k-g(v)\ind_{\{s>0\}},\qquad
 \bigcup_{r\in\mathcal M(v)\setminus B}(N^+(r)\cap U(v))\subseteq L.
\end{equation}
Every exceptional partner $r\in B$ admits the convenient direction
$r\to v$. Residual zero gives $t(v)=0$, residual one gives $t(v)\le1$,
and $\eta(v)=t(v)=1$ gives $P(v)\ne\varnothing$ and
$N^{--}(v)=\partial P(v)$.
\end{lemma}
\begin{proof}
The gap bound follows by deleting one outgoing arc: other first sets are
unchanged and their second sets can only shrink, while the deleted tail's
gap decreases by at most two. The external-boundary deletion argument is
Lemma~\ref{S-lem:boundary}; the original target partition yields
Theorem~\ref{S-thm:main} and the residual identity. The shared remainder
estimate is Lemma~\ref{S-lem:partner-head-budget}; it subtracts the original
second degree if $s=0$, and the second degree forced by minimality after the
nonempty deletion if $s>0$. Its union bound counts distinct heads.
Zero-residual protection and the unit bound are proved in the supplementary
local-budget section. At residual one, a nonempty deletion gives $t=0$;
otherwise equality forces $J$ empty. The empty-$P$ identity gives a whole
vertex and $t=0$, proving the final assertion when $t=1$.
\end{proof}

The set $L$ concerns an original deletion. It is never recomputed in a
closure graph and then assigned the original minimality hypothesis.
In particular, an exceptional missing far head can lie in $X$; no inclusion
$X\subseteq N^-(v)$ is assumed.

\subsection{A residual-four threshold for zero completed gap}
Let $S_v=\mathcal M(v)\cap U(v)$. For nonempty $P(v)$,
\begin{equation}\label{eq:paper-missing-far}
 |S_v|\le k+|L|\le\eta(v)-k-g(v)\ind_{\{s>0\}}.
\end{equation}
Indeed, a nonexceptional missing far head belongs to $E$ and cannot become
reachable when arcs are deleted, so it lies in $L$. When $k=0$, in fact
$\mathcal F(v)\subseteq L$. A far incoming head $y$ that is not protected
has a predecessor $r$ not pointing to $v$. It cannot be an outneighbour of
$v$, since that would reach $y$ in two steps. Thus it is a missing partner,
and the union bound in \eqref{eq:paper-packet-bound} applies.

\begin{theorem}[Zero-gap residual threshold]\label{thm:paper-eta-four}
At a zero-gap feed $f$ of the weak-closure completion, let
$m=|\mathcal M_Q(f)|$. Then
\[
 \eta_D(f)\ge m+1\ge4.
\]
If $\eta_D(f)=4$, its original packet necessarily satisfies
\[
 P_D(f)\ne\varnothing,\quad m=3,\quad k=s=0,\quad g_D(f)=1,
 \quad t_D(f)=4.
\]
Furthermore $|P_D(f)|=2\mu_D(f)-4\ge2$,
$N_D^{--}(f)=\partial P_D(f)$, $L=\mathcal F_D(f)$,
and $|\partial^2P_D(f)|=|\partial P_D(f)|-1$.
\end{theorem}
\begin{proof}
Theorem~\ref{thm:paper-original-cycle} gives $m\ge3$,
$\mathcal M_Q(f)\subseteq S_f$, and $t_D(f)\ge g_D(f)+m$.
If $P(f)$ is empty, the residual identity gives
$\eta(f)=2\mu(f)+g(f)-1\ge2m+g(f)-1\ge m+1$.
For nonempty $P(f)$ and $k>0$, \eqref{eq:paper-missing-far} gives
$\eta(f)\ge m+k+g(f)\ind_{\{s>0\}}\ge m+1$.
For $k=0$, the inclusion $\mathcal F(f)\subseteq L$ gives
$\eta(f)\ge g(f)+m+g(f)\ind_{\{s>0\}}\ge m+1$.

At residual four these bounds force $P(f)\ne\varnothing$, $m=3$,
$s=0$ and $k\le1$. If $k=1$, all three original missing far heads are
exactly the terminal partners, the unique exceptional partner is among
them, and the other two already consume the bound $|L|\le2$.
Their original induced graph has a directed cycle, hence a triangle.
The exceptional vertex receives an original arc from a nonexceptional
partner on this triangle. The shared head bound puts it in $L$ too,
a contradiction. Thus $k=0$; now $g+m\le t\le4$ gives $g=1,t=4$.
The residual identity gives the stated size of $P$. The partition gives
$N^{--}=\partial P$, and $\mathcal F\subseteq L$, $|L|\le4$ gives
$L=\mathcal F$. Here $\widehat D=D$ and
$|E|=d^{++}(f)+4=\mu(f)+|\partial^2P|$; combine this with
$\eta=\mu-1+|\partial P|-d^{++}=4$ to obtain the final equality.
\end{proof}

This threshold does not apply to strictly negative completed gaps. At strong
Good equality the supplementary original path theorem gives a finer statement:
if $\eta(f)=2$, every terminal partner is an original nongood second vertex;
if $\eta(f)=3$, there is at most one terminal far partner and it receives an
original arc from such a second partner. The unique far partner at residual
three may be exceptional. These conclusions leave both the all-second case
and the strictly positive Good-slack case unresolved.

\subsection{Actual failure sources and the residual margin}
For a missing direction $a\to b$, its original failure set is
$\{x:x\to a,\ b\in U_D(x)\}$. Let $\mathcal A$ be the vertices occurring
in the failure set of a direction on a bad missing pair. Opposite directional
witnesses are nonadjacent. The cross-witness graph joins every such pair of
witnesses. A missing pair is flexible if it is bad or convenient in both
directions.

\begin{theorem}[One uncontrolled actual source]\label{thm:paper-one-exception}
In an arbitrary oriented graph, suppose some vertex $u$ has the following
properties: every actual source in $\mathcal A\setminus\{u\}$ has $t\le1$,
and every cross-witness pair between two such sources is flexible. Then the
graph has a Seymour vertex. In MIN, at least two distinct actual failure
sources have residual at least two. In particular every positive residual
allocation $(m,1^k)$ is excluded for $m\ge2$ and finite $k$.
\end{theorem}
\begin{proof}
The complete argument is Theorem~\ref{S-thm:one-exception-short-general}.
Its first optimization uses all canonical protection-reference completions,
maximizing the forward score and then the position of $u$. A unit-budget
feed is balanced. An opposite witness other than $u$ makes a flexible
incoming arc backward after sedimentation, whereas witness $u$ moves $u$
strictly later. Thus the selected feed must be $u$, which therefore has
no original protection successor. A separate fixed completion avoiding
bad failures at $u$, followed by a median maximizing its position, yields
the second contradiction. The fixed-cover argument is proved independently
in Theorem~\ref{S-thm:protected-independent-heavy-cover}.

In MIN, an actual source has positive residual; residual one has $t\le1$,
and the unit--unit opposite-witness pair is bad. If at most one actual
source had residual at least two, the general result would apply. The
allocation consequence follows even when inactive residuals are unrestricted.
\end{proof}

\begin{theorem}[Parity-sensitive missing-pair margin]\label{thm:paper-residual-seven}
Every MIN graph satisfies $R=2M-n\ge7$, with $R\ge8$ at even order.
Every arc-set-minimal all-deficient graph, whether strong or not, satisfies
\[
 M\ge\left\lceil n/2\right\rceil+3+\ind_{\{n\text{ even}\}}.
\]
If $o$ vertices lie outside a sink strong component, the right side can be
increased by $o$. Equality in the uniform bound implies strong connectivity.
\end{theorem}
\begin{proof}
The supplementary two-support theorem gives at least three positive
residual vertices (Corollary~\ref{S-cor:three-positive-residual-support}).
Theorem~\ref{thm:paper-one-exception} requires two residuals at least two.
Hence, at total residual at most six, the only possible positive partitions are
\[
 (2,2,1),\quad(3,2,1),\quad(2,2,2),\quad(2,2,1,1).
\]
The first two are excluded by the general $(m,2,1)$ theorem
(Theorem~\ref{S-thm:no-three-two-one}); the last two by
Theorems~\ref{S-thm:eta222-excluded} and \ref{S-thm:no-two-two-one-one}.
These are complete graph arguments, not an enumeration of possible orders.
Thus $R\ge7$; parity gives $R\ge8$ for even $n$.

For a sink component of order $k$ with $o=n-k$, the inherited minimality
and the outside missing-pair count in Theorem~\ref{S-thm:sink} give
\[
 M(D)\ge\left\lceil k/2\right\rceil+3+\ind_{\{k\text{ even}\}}
               +\left\lceil o(k-2)/2\right\rceil.
\]
An all-deficient graph has minimum outdegree at least two, so $k\ge5$.
Using $k-2\ge3$ and checking the two parities yields at least
$\lceil n/2\rceil+o+3+\ind_{\{n\text{ even}\}}$.
If $o>0$ this is strictly stronger than the uniform bound.
\end{proof}

Thus at $n=2\delta+3$ one has $M\ge\delta+5$, and at $n=2\delta+4$
one has $M\ge\delta+6$. The remaining residual-seven partitions after
these filters are
\[
 (3,3,1),\ (3,2,2),\ (3,2,1,1),\ (2,2,2,1),\ (2,2,1,1,1).
\]
Their weak-closure median feeds must have strictly negative completed gap,
by Theorem~\ref{thm:paper-eta-four}. None of these five partitions is
excluded by that observation alone.

\subsection{An additional global constraint on a near-whole feed}
Suppose a weak-closure feed $f$ has one original missing partner $r$.
Terminal whole-feed exclusion leaves $r$ still missing. Absence of a weak
$r\to f$ step supplies a current predecessor $p$ of $r$ not pointing to
$f$. Since $f$ is adjacent to every other vertex, $f\to p$; exact feed
neighbourhood invariance gives $r\in N_D^{++}(f)$.
If also $P_D(f)$ is empty, each unprotected original far head $y$ must
receive an original arc from $r$, the only possible missing-partner witness
against protection. Put $O=N_D^+(f)$, $d=d_D^+(f)$, $g=g_D(f)$ and $t=t_D(f)$.
Then $\mathcal F_D(f)\subseteq\partial_D^2O$. Every member of $O$ reaches
$f$ originally within two steps, so $f\in\partial_D^2O$ as well. Therefore
\[
 t+1\le|\partial_D^2O|<|\partial_D O|=d-g,
 \qquad d\ge g+t+2,\qquad n\ge g+3t+5.
\]
This uses the global original boundary inequality. It is stronger than
the constraints obtained by deleting only arcs with tail $f$, but still
leaves unbounded parameter ranges.
